% =====================================================================
% pertemuan-13.tex
% Pertemuan 13 - Analisis Semantik: Attribute Grammar, SDD, Translation Scheme
% Sub-CPMK 1.5 - Praktik terbimbing + kuis konsep
% =====================================================================
\documentclass[aspectratio=169,10pt,t]{beamer}
\usepackage{config/beamerthemeubb}

\UBBmeeting{13}

\begin{document}

\UBBtitleframe{13}{Analisis Semantik: Attribute Grammar, SDD, dan Translation Scheme}{Praktik terbimbing \textperiodcentered{} Kuis konsep}

% ---------------------------------------------------------------------
\begin{frame}{Capaian Pertemuan}
  \textbf{CPMK induk:} CPMK 01 --- Menjelaskan tahapan kompilasi
  (\textit{lexical}, \textit{syntax}, \textit{semantic analysis}).

  \vspace{3pt}
  \begin{block}{Sub-CPMK yang dinilai pada pertemuan ini}
    \small
    \textbf{Sub-CPMK 1.5:} Membuat tabel aturan semantik berbasis
    \textit{syntax-directed definition} untuk ekspresi dan pernyataan
    subset bahasa C, lalu menguji keluarannya pada pohon sintaks.
  \end{block}
\end{frame}

% ---------------------------------------------------------------------
\begin{frame}{Agenda Pertemuan (150 menit)}
  \small
  \begin{tabular}{@{}p{1.4cm}p{12.3cm}@{}}
    \toprule
    \textbf{Menit} & \textbf{Kegiatan} \\
    \midrule
    25 & Ceramah: ruang lingkup semantik dan attribute grammar \\
    30 & Latihan: atribut \textit{synthesized} dan \textit{inherited} \\
    35 & Lokakarya: tabel SDD ekspresi dan pernyataan \\
    25 & Latihan: menguji aturan pada pohon sintaks \\
    20 & Kuis konsep dan pembahasan \\
    15 & Kode antara: postfix dan three-address, lalu penutup \\
    \bottomrule
  \end{tabular}

  \vspace{6pt}
  {\scriptsize\color{ubbblue!75}Porsi latihan dan praktik 90 menit;
  kuis konsep 20 menit. Keluaran berupa tabel aturan dan pohon
  teranotasi.}
\end{frame}

% ---------------------------------------------------------------------
\begin{frame}{Posisi Pertemuan pada Rantai Proyek}
  \UBBchain{4}

  \vspace{2pt}
  \begin{itemize}
    \item Pertemuan ini membangun simpul \textbf{semantik}: memeriksa
          makna program di atas AST.
    \item Tabel aturan semantik menjadi dasar implementasi
          \textit{type checker} pada Pertemuan 14.
    \item Keluaran berupa tipe dan error yang dipropagasikan.
  \end{itemize}
\end{frame}

% ---------------------------------------------------------------------
\begin{frame}{Tugas Analisis Semantik}
  \scriptsize
  \setlength{\tabcolsep}{4pt}
  \begin{tabular}{@{}p{4.0cm}p{9.2cm}@{}}
    \toprule
    \textbf{Tugas} & \textbf{Pertanyaan yang dijawab} \\
    \midrule
    Resolusi nama & Apakah identifier sudah dideklarasikan? \\
    Pemeriksaan tipe & Apakah operand memiliki tipe yang kompatibel? \\
    Pemeriksaan fungsi & Apakah jumlah dan tipe argumen sesuai? \\
    Kontrol alur & Apakah \texttt{break} berada di dalam loop? \\
    Inisialisasi & Apakah variabel dipakai sebelum diberi nilai? \\
    Keterjangkauan & Apakah semua jalur fungsi mengembalikan nilai? \\
    \bottomrule
  \end{tabular}

  \vspace{5pt}
  {\scriptsize\color{ubbblue!75}Semua tugas ini berada di luar daya
  ekspresif CFG; diperlukan tabel simbol dan aturan tipe.}
\end{frame}

% ---------------------------------------------------------------------
\begin{frame}{Informasi dan Urutan Analisis}
  \begin{columns}[T,onlytextwidth]
    \begin{column}{0.46\textwidth}
      \textbf{Informasi yang dibutuhkan}
      \begin{itemize}
        \item AST atau pohon urai;
        \item tabel simbol dengan scope dan tipe;
        \item aturan tipe setiap operator;
        \item mekanisme pelaporan kesalahan.
      \end{itemize}
    \end{column}
    \begin{column}{0.50\textwidth}
      \textbf{Urutan analisis}
      \begin{enumerate}
        \item bangun tabel simbol dan resolusi nama;
        \item tentukan tipe setiap ekspresi;
        \item periksa kecocokan tipe;
        \item periksa kontrol alur;
        \item laporkan kesalahan dan lanjutkan.
      \end{enumerate}
    \end{column}
  \end{columns}

  \vspace{4pt}
  {\footnotesize\color{ubbblue!80}Analisis tidak menghasilkan kode;
  tujuannya memastikan program valid dan mengumpulkan informasi untuk
  fase berikutnya.}
\end{frame}

% ---------------------------------------------------------------------
\begin{frame}{\textit{Attribute Grammar}}
  \begin{block}{Definisi}
    \small
    \textit{Attribute grammar} adalah tata bahasa bebas konteks yang
    setiap simbolnya memiliki atribut dan setiap produksinya memiliki
    aturan semantik untuk menghitung atribut tersebut.
  \end{block}

  \begin{itemize}
    \item Atribut dapat berupa tipe, nilai, alamat, kode, atau
          informasi lain.
    \item Setiap atribut memiliki domain nilai dan aturan evaluasi.
    \item Ketergantungan antar-atribut membentuk graf berarah.
    \item Evaluasi mengikuti urutan topologis graf; siklus harus
          dihindari.
  \end{itemize}
\end{frame}

% ---------------------------------------------------------------------
\begin{frame}{Atribut \textit{Synthesized} dan \textit{Inherited}}
  \begin{columns}[T,onlytextwidth]
    \begin{column}{0.48\textwidth}
      \begin{block}{\textit{Synthesized}}
        \footnotesize
        Dihitung dari atribut anak. Nilai mengalir dari daun ke akar.\\
        Contoh: \(E.type=\text{rule}(E_1.type,E_2.type)\).
      \end{block}
    \end{column}
    \begin{column}{0.48\textwidth}
      \begin{block}{\textit{Inherited}}
        \footnotesize
        Dihitung dari orang tua atau saudara. Nilai mengalir dari akar
        ke daun atau antar-saudara.\\
        Contoh: \(\text{stmt}.env=\text{block}.env\).
      \end{block}
    \end{column}
  \end{columns}

  \vspace{4pt}
  \begin{center}
    \begin{tikzpicture}[
        nd/.style={draw=ubbblue!40,fill=ubbblue!7,rounded corners=2pt,
          minimum width=1.8cm,minimum height=0.72cm,align=center,
          font=\scriptsize},
        ar/.style={-{Stealth[length=1.5mm]},draw=ubbblue!55,thick}
      ]
      \node[nd] (root) at (0,0) {\(E\)\\type = float};
      \node[nd] (l) at (-2.2,-1.1) {\(E_1\)\\type = int};
      \node[nd] (r) at (2.2,-1.1) {\(E_2\)\\type = float};
      \draw[ar] (l) -- node[left]{\scriptsize synthesized} (root);
      \draw[ar] (r) -- node[right]{\scriptsize synthesized} (root);
    \end{tikzpicture}
  \end{center}
\end{frame}

% ---------------------------------------------------------------------
\begin{frame}{\textit{Syntax-Directed Definition} (SDD)}
  \begin{block}{Definisi}
    \small
    SDD adalah tata bahasa dengan atribut dan aturan semantik. Setiap
    produksi \(A \to \alpha\) memiliki aturan untuk menghitung atribut
    simbol-simbolnya.
  \end{block}

  \begin{itemize}
    \item SDD mendeskripsikan \emph{apa} yang dihitung, bukan
          \emph{kapan} perhitungan dilakukan.
    \item Urutan evaluasi ditentukan oleh ketergantungan antar-atribut.
    \item SDD S-attributed hanya memakai atribut synthesized;
          dapat dievaluasi post-order.
    \item SDD L-attributed memakai inherited yang bergantung pada
          kiri atau orang tua; dievaluasi kiri-ke-kanan.
  \end{itemize}
\end{frame}

% ---------------------------------------------------------------------
\begin{frame}{Tabel SDD untuk Ekspresi (Artefak Sub-CPMK 1.5)}
  \scriptsize
  \setlength{\tabcolsep}{4pt}
  \begin{tabular}{@{}p{4.2cm}p{9.0cm}@{}}
    \toprule
    \textbf{Produksi} & \textbf{Aturan semantik} \\
    \midrule
    \(E \to E_1+E_2\) &
      \(E.type=\text{arith}(E_1.type,E_2.type)\) \\
    \(E \to E_1*E_2\) &
      \(E.type=\text{arith}(E_1.type,E_2.type)\) \\
    \(E \to \code{id}\) &
      \(E.type=\text{lookup}(\code{id}.name)\) \\
    \(E \to \code{INT\_LIT}\) & \(E.type=\text{int}\) \\
    \(E \to \code{FLOAT\_LIT}\) & \(E.type=\text{float}\) \\
    \(E \to (E_1)\) & \(E.type=E_1.type\) \\
    \bottomrule
  \end{tabular}

  \vspace{5pt}
  {\scriptsize\color{ubbblue!75}\texttt{arith} mengembalikan
  \texttt{float} jika salah satu operand float, \texttt{int} jika
  keduanya int/char, dan \texttt{error} jika operand tidak numerik.}
\end{frame}

% ---------------------------------------------------------------------
\begin{frame}{Tabel SDD untuk Pernyataan}
  \scriptsize
  \setlength{\tabcolsep}{4pt}
  \begin{tabular}{@{}p{3.4cm}p{6.2cm}p{3.6cm}@{}}
    \toprule
    \textbf{Produksi} & \textbf{Aturan semantik} &
      \textbf{Keluaran uji} \\
    \midrule
    \(S \to \code{id}=E\) &
      \(S.type=\text{check\_assign}(\text{lookup}(\code{id}),E.type)\) &
      tipe \(S\) atau error \\
    \(S \to \code{if}(E)\ S_1\) &
      \(\text{check\_bool}(E.type);\ S.type=S_1.type\) & tipe \(S\) \\
    \(S \to \code{while}(E)\ S_1\) &
      \(\text{check\_bool}(E.type);\ S.type=S_1.type\) & tipe \(S\) \\
    \(S \to \{S_1\ S_2\}\) &
      \(S.type=\text{seq}(S_1.type,S_2.type)\) & tipe \(S\) \\
    \(D\) deklarasi &
      \(\text{declare}(\code{id},\code{type})\) & entri tabel simbol \\
    \bottomrule
  \end{tabular}

  \vspace{5pt}
  {\scriptsize\color{ubbblue!75}Tabel ini melengkapi tabel ekspresi;
  keduanya menjadi artefak Sub-CPMK 1.5.}
\end{frame}

% ---------------------------------------------------------------------
\begin{frame}[fragile]{Menguji Aturan pada Pohon Sintaks}
  \begin{columns}[T,onlytextwidth]
    \begin{column}{0.42\textwidth}
      \begin{lstlisting}[language=C,numbers=none,
        caption={Program uji aturan semantik}]
int a = 2;
float b = 3.5;
a = b + 1;
      \end{lstlisting}
    \end{column}
    \begin{column}{0.55\textwidth}
      \scriptsize
      \setlength{\tabcolsep}{3pt}
      \begin{tabular}{@{}p{2.1cm}p{3.3cm}p{1.8cm}@{}}
        \toprule
        \textbf{Simpul} & \textbf{Aturan} & \textbf{Hasil} \\
        \midrule
        Deklarasi \(a\) & declare\((a,\code{int})\) & \(a:\code{int}\) \\
        Deklarasi \(b\) & declare\((b,\code{float})\) & \(b:\code{float}\) \\
        \(b+1\) & arith(float, int) & float \\
        \(a=b+1\) & check\_assign(int, float) & error \\
        \bottomrule
      \end{tabular}
    \end{column}
  \end{columns}

  \vspace{4pt}
  {\footnotesize\color{ubbblue!80}Setiap baris merujuk simpul pohon
  konkret, bukan hanya nama aturan.}
\end{frame}

% ---------------------------------------------------------------------
\begin{frame}{S-Attributed dan L-Attributed}
  \small
  \setlength{\tabcolsep}{4pt}
  \begin{tabular}{@{}p{3.0cm}p{4.8cm}p{4.8cm}@{}}
    \toprule
    \textbf{Aspek} & \textbf{S-attributed} & \textbf{L-attributed} \\
    \midrule
    Atribut & Synthesized & Synthesized dan inherited \\
    Evaluasi & Post-order & Kiri-ke-kanan \\
    Implementasi & Mudah & Perlu konteks \\
    Contoh & Tipe ekspresi & Lingkungan scope \\
    \bottomrule
  \end{tabular}

  \vspace{5pt}
  \begin{block}{Evaluasi pada proyek Subset-C}
    \footnotesize
    Evaluasi dilakukan setelah AST terbentuk, dengan traversal.
    Pendekatan ini memisahkan parsing dari analisis dan memudahkan
    pengujian.
  \end{block}
\end{frame}

% ---------------------------------------------------------------------
\begin{frame}{\textit{Translation Scheme}}
  \begin{block}{Definisi}
    \small
    \textit{Translation scheme} adalah tata bahasa dengan aksi
    semantik yang ditempatkan pada posisi tertentu dalam produksi.
    Aksi dijalankan ketika posisi tersebut dicapai selama parsing.
  \end{block}

  \begin{exampleblock}{Contoh}
    \footnotesize
    \(E \to E_1+\{\,E.type=\text{arith}(E_1.type,E_2.type);\,\}E_2\).\\
    Aksi dijalankan setelah \(E_1\) dan sebelum \(E_2\) diproses.
  \end{exampleblock}

  \vspace{2pt}
  {\footnotesize SDD mendeskripsikan aturan; translation scheme
  mendeskripsikan urutan eksekusi.}
\end{frame}

% ---------------------------------------------------------------------
\begin{frame}[fragile]{Aksi Semantik pada Bison}
  \begin{columns}[T,onlytextwidth]
    \begin{column}{0.52\textwidth}
      \begin{lstlisting}[language=Bison,numbers=none,
        caption={Translation scheme pada Bison}]
expr : expr T_PLUS expr
     {
         $$ = ast_new(AST_BINARY, $2.token);
         $$->left = $1;
         $$->right = $3;
     }
     ;
      \end{lstlisting}
    \end{column}
    \begin{column}{0.44\textwidth}
      \begin{itemize}
        \item Aksi dijalankan saat aturan direduksi.
        \item Kedua operand sudah tersedia.
        \item Pola S-attributed yang khas.
        \item Logika kompleks dipindahkan ke fungsi bantu.
        \item Grammar tetap ringkas dan dapat diuji.
      \end{itemize}
    \end{column}
  \end{columns}
\end{frame}

% ---------------------------------------------------------------------
\begin{frame}[fragile]{Modul Tipe \file{types.h}}
  \begin{columns}[T,onlytextwidth]
    \begin{column}{0.52\textwidth}
      \lstinputlisting[language=C,firstline=7,lastline=21,
        firstnumber=7,
        caption={Enum dan struktur tipe pada \file{types.h}},
        label={lst:p13-types}]{../buku_ajar/code/subsetc/include/types.h}
    \end{column}
    \begin{column}{0.44\textwidth}
      \begin{itemize}
        \item \texttt{TYPE\_ERROR} sebagai nilai propagasi.
        \item Jika satu operand error, hasil juga error.
        \item Kesalahan tidak berlipat ganda.
        \item \texttt{type\_arith} dan \texttt{type\_assignable}
              mengimplementasikan aturan.
      \end{itemize}
    \end{column}
  \end{columns}
\end{frame}

% ---------------------------------------------------------------------
\begin{frame}{Notasi Postfix}
  \small
  \begin{itemize}
    \item Operator ditulis setelah operand.
    \item \((a+b)*(c+d)\) menjadi \texttt{a b + c d + *}.
    \item Evaluasi memakai stack: operand didorong; operator mengambil
          dua operand teratas, lalu mendorong hasilnya.
  \end{itemize}

  \vspace{3pt}
  \begin{tabular}{@{}l l l@{}}
    \toprule
    \textbf{Infix} & \textbf{Postfix} & \textbf{Catatan} \\
    \midrule
    \texttt{a + b} & \texttt{a b +} & satu operator \\
    \texttt{a + b * c} & \texttt{a b c * +} & perkalian dahulu \\
    \texttt{(a + b) * c} & \texttt{a b + c *} & kurung mengubah urutan \\
    \texttt{a = b + c * d} & \texttt{b c d * + a =} & assignment \\
    \bottomrule
  \end{tabular}

  \vspace{3pt}
  {\footnotesize\color{ubbblue!80}Postfix tidak memerlukan tanda
  kurung karena urutan operasi ditentukan posisi operator.}
\end{frame}

% ---------------------------------------------------------------------
\begin{frame}{N-Tuple dan Three-Address Code}
  \begin{columns}[T,onlytextwidth]
    \begin{column}{0.52\textwidth}
      \textbf{Quadruple untuk \texttt{a = b + c * d}}
      \scriptsize
      \setlength{\tabcolsep}{3pt}
      \begin{tabular}{@{}c l l l l@{}}
        \toprule
        \textbf{No} & \textbf{Op} & \textbf{O1} & \textbf{O2} &
        \textbf{Hasil} \\
        \midrule
        1 & \texttt{*} & \texttt{c} & \texttt{d} & \texttt{t1} \\
        2 & \texttt{+} & \texttt{b} & \texttt{t1} & \texttt{t2} \\
        3 & \texttt{=} & \texttt{t2} & --- & \texttt{a} \\
        \bottomrule
      \end{tabular}
    \end{column}
    \begin{column}{0.44\textwidth}
      \textbf{Three-address code}
      \begin{block}{}
        \footnotesize
        \texttt{t1 = c * d}\\
        \texttt{t2 = b + t1}\\
        \texttt{a = t2}
      \end{block}
    \end{column}
  \end{columns}

  \vspace{4pt}
  \begin{itemize}
    \item Triple merujuk hasil antara dengan nomor; quadruple menyimpan
          nama hasil.
    \item Three-address code dekat dengan instruksi assembly.
  \end{itemize}
\end{frame}

% ---------------------------------------------------------------------
\begin{frame}{Kuis Konsep}
  \small
  \begin{enumerate}
    \setlength{\itemsep}{3pt}
    \item Apa perbedaan atribut \textit{synthesized} dan
          \textit{inherited}?
    \item Mengapa CFG tidak dapat memeriksa apakah identifier sudah
          dideklarasikan?
    \item Tulis aturan SDD untuk tipe \(E \to E_1+E_2\).
    \item Apa fungsi \texttt{TYPE\_ERROR} dalam propagasi error?
    \item Kapan \textit{translation scheme} lebih praktis daripada SDD?
  \end{enumerate}

  \vspace{4pt}
  {\footnotesize\color{ubbblue!75}Kerjakan mandiri selama 8 menit,
  lalu bahas bersama. Kunci tersedia pada versi handout.}
\end{frame}

\mode<handout>{
\begin{frame}{Kunci Kuis Konsep}
  \footnotesize
  \begin{enumerate}
    \setlength{\itemsep}{2pt}
    \item \textit{Synthesized} dihitung dari anak; \textit{inherited}
          dari orang tua atau saudara.
    \item CFG hanya memeriksa bentuk; deklarasi memerlukan konteks dan
          tabel simbol.
    \item \(E.type=\text{arith}(E_1.type,E_2.type)\).
    \item Meredam pesan turunan: operand error menghasilkan error tanpa
          laporan baru.
    \item Translation scheme lebih dekat ke implementasi dan memuat
          urutan eksekusi aksi.
  \end{enumerate}
\end{frame}
}

% ---------------------------------------------------------------------
\begin{frame}{Kegiatan Mahasiswa di Kelas: P4-A}
  \begin{block}{Tabel aturan semantik dan kode antara}
    \footnotesize
    \begin{enumerate}
      \setlength{\itemsep}{1pt}
      \item Lengkapi tabel aturan ekspresi dan pernyataan.
      \item Pilih program dengan deklarasi, assignment, dan
            percabangan/perulangan; gambar pohon sintaksnya.
      \item Terapkan setiap aturan pada simpul dan catat atribut
            \texttt{type} serta \texttt{code}.
      \item Tuliskan postfix dan three-address code untuk setiap
            ekspresi.
      \item Tuliskan quadruple dan three-address code untuk setiap
            pernyataan.
      \item Uji minimal 15 ekspresi dan 5 pernyataan.
      \item Kumpulkan tabel, pohon teranotasi, dan keluaran kode antara.
    \end{enumerate}
  \end{block}
\end{frame}

% ---------------------------------------------------------------------
\begin{frame}{Tugas dan Artefak}
  \begin{block}{Artefak Sub-CPMK 1.5}
    Tabel aturan semantik berbasis SDD untuk ekspresi dan pernyataan,
    pohon sintaks teranotasi, dan keluaran pengujian aturan.
  \end{block}

  \small
  \begin{itemize}
    \item Tabel ekspresi dan pernyataan lengkap.
    \item Pohon sintaks teranotasi dengan tipe dan atribut kode.
    \item Minimal 15 ekspresi dan 5 pernyataan diuji.
    \item Postfix, quadruple, dan three-address code tersedia.
    \item Laporan dan keluaran pengujian diserahkan.
  \end{itemize}

  \vspace{3pt}
  {\footnotesize\color{ubbblue!75}P4-A menjadi bagian Praktikum 4
  (12\%) yang dilaporkan pada Pertemuan 14.}
\end{frame}

% ---------------------------------------------------------------------
\begin{frame}{Rangkuman dan Jembatan}
  \begin{itemize}
    \item Analisis semantik memeriksa makna: deklarasi, tipe, fungsi,
          kontrol alur, dan inisialisasi.
    \item Attribute grammar memakai atribut synthesized dan inherited.
    \item SDD mendeskripsikan aturan; translation scheme menentukan
          urutan eksekusi.
    \item Tabel aturan ekspresi dan pernyataan diuji pada pohon
          sintaks konkret.
    \item \texttt{TYPE\_ERROR} meredam pesan turunan dan menjaga
          diagnostik tetap fokus.
  \end{itemize}

  \vspace{4pt}
  \begin{exampleblock}{Pertemuan berikutnya}
    Pertemuan 14 mendalami symbol table bersarang, manajemen scope,
    dan type checking dalam C, lalu menutup Praktikum 4.
  \end{exampleblock}
\end{frame}

\end{document}
